
\section{SWEEP Rule Decoder on the Toric, version III}
Consider the presentation of a $d$-dimensional torus by the Cayley graph $G = (V,E)$ of the $\left( \mathbb{Z}_{n}^{d}, + \right)$ group generated by the elements $S = \{ 1_{i} \colon i \in [d] \}$ (where $1_{i}$ is the vector consisting of one at the $i$th coordinate and zero elsewhere). We say that an edge $\{x, y\}$ is positive relative to a plaquette $\{v, g_{1}v, g_{2}g_{1}v, g_{2}v\}$, defined by a rooted vertex $v \in \mathbb{Z}_{n}^{d}$ and the generators $g_{1}, g_{2} \in S$, if $x = v$ and $y$ is either $g_{1}v$ or $g_{2}v$. 

In our code, bits are set on the plaquette, and the edges correspond to checks that restrict the parity of the bits on their support to be even. Denote the set of plaquettes by $\mathcal{F}$, and let us abuse the notation by using it to refer also to the bits.


\begin{definition}[$v$ adds $u$ at time $t$]
  Let $D^{(t)} \subset \mathcal{F}$ be the bits that deviate from the ideal running at time $t$ \ctt{Define it!}. Let $\Gamma^{(t)} \subset V$ be the vertices that support the plaquettes in $D^{(t)}$. Namely, $\Gamma^{(t)} = \{ v \colon \exists f \in D^{(t)} \text{ s.t. } v \in f\}$. In other words, $\Gamma^{(t)}$ are the vertices that, at time $t$, are adjacent to a deviated bit. Similarly, denote by $\Gamma^{* (t)}$ the vertices that at time $t$ adjacent to deviate bits which are not belongs to the faults, i.e $\Gamma^{*(t)} = \{ v : \exists f \in D^{(t)} / \mathcal{E}^{(t)}\text{ s.t. } v \in f\}$
We say that for vertices $v, u \in V$, $v$ adds $u$ at time $t$, if the following conditions hold:
\begin{enumerate}
  \item $v \in \Gamma^{t-1}$  and $ u \in \Gamma^{* (t)}$
  \item There exists $f \in \mathcal{F}$ such that  $v, u \in f$.  
  \item $v \le u $, namely $v_{i} \le u_{i}$ for any $i\in [d]$.
\end{enumerate}
For a time $t$, we define the function $P_{t} \colon  V \rightarrow 2^{V}$ that, for a vertex $u$, returns the set of vertices that add $u$.
\end{definition}

\begin{claim}[Shrinking Lemma]
  For any $u \in \Gamma^{(*) t}$, there are $v_{1}, \ldots, v_{d+1} \in P_{t}(u)$, not necessarily different, such that for any $i \le d$, it holds that $L_{i}(v_{i}) = L_{i}(u)$, and for $i = d+1$, it holds that $L_{d+1}(v_{d+1}) \ge L_{d+1}(u) + 1$.
\end{claim}

To prove it, we will need the following claim:
\begin{claim}
  \label{claim:inprev}
For any $u \in \Gamma^{* (t)}$, there is at least one $v \in P_{t}(u)$ such that $v < u$.
\end{claim}
\begin{proof}
  Since $u \in \Gamma^{*(t)}$ there must be $f \in D^{(t)}/\mathcal{E}^{(t)}$ such $u \in f$.
  \begin{itemize}
    \item If $f \in D^{(t-1)}$, then it least one of its positive edge
  \end{itemize}
\end{proof}

\begin{proof}
  Consider $u \in \Gamma^{* (t)}$. By \cref{claim:inprev} there is at least one $v \in P_{t}(u)$ satisfies $v < u$, choose it to be $v_{d+1}$. Recall that since the vertices are the points of the d-dimensional lattice embedded in $\mathbb{R}^{d}$, the inequality $v <u$ implies that there exists at least single $i \in[d]$ such $ v+ \mathbf{1}_{i} \le u$. Therefore we get: 
  \begin{equation*}
    \begin{split}
      L_{d+1}(v_{d+1}) = -\braket{ \mathbf{1}, v_{d+1} } = -\sum_{j \in [d]}v_{j} =  -\sum_{j \in [d]/i}v_{j} - v_{i} \ge  -\braket{ \mathbf{1}, u } + 1 = L_{d+1}(u) + 1
    \end{split}
  \end{equation*}
  Now let's split to cases upon whether $u \in \Gamma^{(t-1)}$. 
  \begin{itemize}
    \item If $u \in \Gamma^{(t-1)}$ then picking it for the vertices $v_{1},..,v_{d}$ finishes the proof.
    \item Else, there has to be a plaquette $f \in D^{(t)}/D^{(t-1)}$ such that $v \in f$. Thus, $f$ was flipped at time $t$. Since, by the definition of Toom's rule, a plaquette is flipped only if two of its edges are unsatisfied\footnote{Here the argument works no matter which edges have been flipped.}, we get that at least three vertices of $f$ are adjacent to unsatisfied checks at time $t-1$, namely adjacent to plaquettes in $D^{(t-1)}$\footnote{Those plaquettes are adjacent to $f$.}.
      
Since those vertices are the vertices of $f / \{u\}$, two of them, let's say $a$ and $b$, differ from $u$ at only one coordinate, let's say $j_{a}$ and $j_{b}$, respectively. Therefore, define:
      \begin{equation*}
        \begin{split} v_{i} = 
          \begin{cases}
            v_{a} & a \neq j_{a} \\ 
            v_{b} & \text{else}
          \end{cases}
        \end{split}
      \end{equation*}
  \end{itemize}
  Clearly, we got the equalities for $i \le d$.
\end{proof}

e generalize Toom's rule in the following manner: Each plaquette asks if both of its positive edges point to a syndrome, namely, are unsatisfied. If so, then the bit flips itself; for convenience, we will think of the vertex at the end of the positive edges (in the 2D bottom-left corner) as the vertex that makes the decision and flips, or does not flip, the plaquette. For a vertex $ v \in V $ and plaquette $ f \in \mathcal{F} $, we say that $ v \in f $ if and only if $ v $ is one of the four vertices on $ f $'s corners.
